\subsection{Substrate Query Scaling}
\label{sec:substrate_query_scaling}

\paragraph{Task Description.}
This experiment loads a 400-sample synthetic corpus (128\,B each,
50\,KB) and issues $N = 50$ prefix queries in corpus-position order.
We instrument the prediction machinery: beam confidence (top
continuation log-score), score spread (gap between best and worst
hypothesis — measures diversity), rejection rate (fraction of trie
origins pruned at node-level beam), and origin diversity (how many
tries contribute surviving hypotheses). These metrics expose whether
the substrate's internal state scales gracefully or degrades as the
corpus grows.

\paragraph{Results.}
Beam confidence rises across the query sequence (top-score
$$\mu = -3.087$, $\sigma = 4.826$, median $= -3.466$, range $[-32.716,\,0.000]$, $N = 50$$). Score spread: $$\mu = 0.012$, $\sigma = 0.088$, median $= 0.000$, range $[0.000,\,0.625]$, $N = 50$$. Rejection rate: $$\mu = 0.018$, $\sigma = 0.126$, median $= 0.000$, range $[0.000,\,0.889]$, $N = 50$$.
Origin diversity: $\mu = 1.0$ tries contributing per query.

Figure~\ref{fig:substrate_query_scaling} shows four scaling
perspectives: beam confidence trajectory (top-left), hypothesis
diversity via score spread (top-right), mesh noise via rejection
rate (bottom-left), and effective mesh utilization (bottom-right).

\paragraph{Assessment.}
Beam confidence: $\mu = -3.087$ ($\sigma = 4.826$). Rejection rate:
$0.018$. Origin diversity: $\mu = 1.0$ tries contributing. The
score spread panel shows whether surviving hypotheses are tightly
clustered (low spread = consensus) or divergent (high spread =
competing interpretations).
